\chapter{Relacionamentos: a parte relacional}
\label{cap:relacionamentos}

Até aqui usamos uma tabela só --- isso é uma planilha, não um banco relacional.
O que muda tudo é \textbf{relacionar} tabelas.

\section{O problema que a chave estrangeira resolve}

Se guardássemos o nome do curso dentro de \comando{membro}, teríamos
``Engenharia de Software'' repetido em dezenas de linhas. Aí alguém digita
``Eng. de Software'', outro digita ``engenharia de software'', e a consulta por
curso quebra. Corrigir o nome exigiria alterar todas as linhas.

A solução: o curso vira uma tabela própria, e \comando{membro} guarda apenas o
\comando{id} do curso.

\begin{lstlisting}[style=base,caption={Duas tabelas ligadas por chave},label={lst:fk}]
   curso                          membro
+----+---------------------+   +----+------------+----------+
| id | nome                |   | id | nome       | curso_id |
+----+---------------------+   +----+------------+----------+
|  1 | Eng. de Software    |<--|  1 | Ana Souza  |        1 |
|  2 | Ciencia da Comp.    |<--|  2 | Bruno Lima |        2 |
+----+---------------------+   |  3 | Carla Dias |        1 |
       ^                       +----+------------+----------+
  chave primaria (PK)                          chave estrangeira (FK)
\end{lstlisting}

\begin{conceito}[PK e FK]
A \textbf{chave primária (PK)} identifica a linha \emph{dentro} da própria
tabela. A \textbf{chave estrangeira (FK)} é uma coluna que \emph{aponta} para a
PK de outra tabela. Com ela declarada, o SGBD passa a \textbf{garantir} que
aquele valor existe do outro lado --- isso se chama \textbf{integridade
referencial}.
\end{conceito}

\section{Declarando a chave estrangeira}

\begin{sqlcode}
ALTER TABLE membro
  ADD CONSTRAINT fk_membro_curso
  FOREIGN KEY (curso_id) REFERENCES curso(id)
  ON DELETE SET NULL
  ON UPDATE CASCADE;
\end{sqlcode}

Com a chave estrangeira ativa, o comando abaixo passa a \textbf{falhar} --- e é
exatamente o que queremos:

\begin{sqlcode}
INSERT INTO membro (nome, email, periodo, data_ingresso, curso_id)
VALUES ('Fake', 'fake@liga.edu.br', 1, '2026-01-01', 999);
-- ERROR 1452: Cannot add or update a child row:
--             a foreign key constraint fails
\end{sqlcode}

\begin{table}[H]
\centering
\caption{Ações de \comando{ON DELETE} e \comando{ON UPDATE}}
\label{tab:fk-acoes}
\footnotesize
\begin{tabularx}{\textwidth}{@{}lX@{}}
\toprule
\textbf{Ação} & \textbf{Efeito ao apagar ou alterar o curso} \\
\midrule
\comando{RESTRICT} (padrão) & Impede a operação enquanto houver membro apontando para ele \\
\comando{SET NULL}          & Os membros ficam com \comando{curso\_id = NULL}; exige coluna que aceite \comando{NULL} \\
\comando{CASCADE}           & Propaga: apagar o curso apagaria os membros. \textbf{Use com muito cuidado} \\
\bottomrule
\end{tabularx}
\end{table}

\section{Cardinalidades}

\begin{itemize}
    \item \textbf{1:N (um para muitos)} --- o caso mais comum. Um curso tem
          vários membros; cada membro tem um curso. \textbf{A chave estrangeira
          fica no lado ``muitos''}, ou seja, em \comando{membro}. É o nosso caso.
    \item \textbf{1:1} --- chave estrangeira com \comando{UNIQUE}. Raro.
    \item \textbf{N:N (muitos para muitos)} --- por exemplo, um membro
          participa de vários projetos e um projeto tem vários membros.
          Resolve-se com uma \textbf{tabela intermediária}
          (\comando{membro\_projeto}) que guarda duas chaves estrangeiras.
          Conceito citado, não praticado nesta capacitação.
\end{itemize}

\section{\comando{JOIN}: consultando as duas tabelas juntas}

\comando{JOIN} combina linhas de tabelas diferentes usando a condição de
ligação declarada no \comando{ON}.

\begin{sqlcode}
-- INNER JOIN: so os membros QUE TEM curso
SELECT m.nome AS membro,
       c.nome AS curso,
       m.periodo
FROM   membro m
INNER  JOIN curso c ON m.curso_id = c.id
ORDER  BY c.nome, m.nome;
\end{sqlcode}

\newpage

\begin{sqlcode}
-- LEFT JOIN: TODOS os membros, tenham curso ou nao.
-- Elisa aparece com curso = NULL.
SELECT m.nome AS membro,
       c.nome AS curso
FROM   membro m
LEFT   JOIN curso c ON m.curso_id = c.id;
\end{sqlcode}

\begin{table}[H]
\centering
\caption{Tipos de \comando{JOIN}}
\label{tab:joins}
\footnotesize
\begin{tabularx}{\textwidth}{@{}lX@{}}
\toprule
\textbf{Tipo} & \textbf{Retorna} \\
\midrule
\comando{INNER JOIN} & Só as linhas com correspondência \textbf{nos dois lados} \\
\comando{LEFT JOIN}  & \textbf{Todas} as da tabela da esquerda, mais as correspondentes da direita (\comando{NULL} se não houver) \\
\comando{RIGHT JOIN} & O espelho do \comando{LEFT}; raro, basta inverter a ordem das tabelas \\
\bottomrule
\end{tabularx}
\end{table}

\begin{dica}[Apelidos e a armadilha do \comando{ON}]
\comando{membro m} e \comando{curso c} são \textbf{apelidos}
(\emph{aliases}). Com duas tabelas em jogo, \comando{m.nome} e \comando{c.nome}
deixam claro de qual tabela é cada coluna; sem isso o MySQL responde
\comando{ERROR 1052: Column 'nome' in field list is ambiguous}.

A condição do \comando{ON} é sempre \textbf{FK = PK}. Esquecer o \comando{ON}
produz um produto cartesiano: cada membro cruzado com cada curso.
\end{dica}

\section{Bônus: contar e agrupar}
\label{sec:agregacao}

\begin{dica}
Este conteúdo entra somente se o cronograma permitir
(Seção~\ref{sec:roteiro}). Caso contrário, fica como material de estudo.
\end{dica}

Funções de agregação resumem várias linhas em um valor: \comando{COUNT},
\comando{SUM}, \comando{AVG}, \comando{MIN} e \comando{MAX}.

\begin{sqlcode}
SELECT COUNT(*) FROM membro;      -- quantos membros existem
SELECT AVG(periodo) FROM membro;  -- periodo medio
SELECT MIN(data_ingresso), MAX(data_ingresso) FROM membro;
\end{sqlcode}

\comando{GROUP BY} aplica a agregação \textbf{por grupo}:

\begin{sqlcode}
SELECT c.nome        AS curso,
       COUNT(m.id)   AS total_membros
FROM   curso c
LEFT   JOIN membro m ON m.curso_id = c.id
GROUP  BY c.id, c.nome
ORDER  BY total_membros DESC;
\end{sqlcode}

\begin{conceito}[\comando{WHERE} e \comando{HAVING}]
\comando{WHERE} filtra \textbf{linhas antes} de agrupar; \comando{HAVING}
filtra \textbf{grupos depois} de agregar. Por exemplo:
\comando{... GROUP BY c.id, c.nome HAVING COUNT(m.id) >= 2;}
\end{conceito}
